\usepackage{tikz}
\usepackage{xcolor}
\usepackage{tabularx}
\usepackage{orcidlink}
\usepackage{lineno}
\usepackage{draftwatermark}
\usepackage{fontspec}
\usepackage{fvextra} % fvextra : package latex qui améliore l'affichage brute de verbatim

\setmainfont{LiberationSans}[
  Path = _extensions/Article/fonts/,
  Extension = .ttf,
  UprightFont = *-Regular,
  BoldFont = *-Bold,
  ItalicFont = *-Italic,
  BoldItalicFont = *-BoldItalic
]

\setsansfont{LiberationSans}[
  Path = _extensions/Article/fonts/,
  Extension = .ttf,
  UprightFont = *-Regular,
  BoldFont = *-Bold,
  ItalicFont = *-Italic,
  BoldItalicFont = *-BoldItalic
]

% Amélioration de l'affichage des blocs de code générés par Pandoc
\RecustomVerbatimEnvironment{Highlighting}{Verbatim}{ %highlighting c'est les chunks dans l'environnement Pandoc et verbatim : moteur d'affichage brute (aucune modif)
  breaklines,  % permet un retour à la ligne dans les chunks
  breaksymbolleft={},
  breaksymbolright={},
  commandchars=\\\{\} % commandchars : pour garder le style des chunks
}
% Source : https://distrib-coffee.ipsl.jussieu.fr/pub/mirrors/ctan/macros/latex/contrib/fvextra/fvextra.pdf  (page 6)